
Semantic entropy (SE) has emerged as an effective approach for hallucination detection in large language models, achieving AUROC values between 0.76 and 0.97 across various benchmarks \citep{Farquhar2024Detecting}. SE measures uncertainty by clustering semantically equivalent responses and computing entropy over these clusters. However, SE computation requires generating 10-20 diverse responses per query and performing pairwise entailment checks via natural language inference models, resulting in substantial computational overhead that limits real-time deployment.

Semantic Entropy Probes (SEPs) address this limitation by training lightweight probes on LLM hidden states to predict SE in a single forward pass \citep{Kossen2024Semantic}. Published results indicate that SEPs achieve AUROC within 2-3\% of full SE on TruthfulQA.

This work attempted to extend the SEP approach with SE-Distilled Probes (SEDP), which incorporate semantic similarity features as auxiliary training signals. The hypothesis was that similarity structure, computed on output text rather than internal representations, would provide model-agnostic regularization enabling cross-model transfer. The initial experiment tested whether SEDP could achieve meaningful SE correlation (Spearman rho >= 0.3) on same-model evaluation as an existence proof.

The experiment produced negative results. SEDP achieved Spearman rho = 0.0843 with true SE (p = 0.283), failing to reach the threshold of 0.3 by a margin of 72\%. Both SEDP and the baseline SEP achieved AUROC approximately 0.52, effectively indistinguishable from random classification. The delta between SEDP and SEP was +0.0009 in correlation and +0.0004 in AUROC, indicating negligible contribution from similarity features.

The 39\% gap between the achieved AUROC (0.52) and published benchmarks (approximately 0.85) represents a substantial discrepancy. This gap may arise from configuration sensitivity---the experiment tested layer 25, Token-Before-Generation (TBG) position, and logistic regression architecture without systematic ablation across alternative configurations.

This work makes the following contributions:

\begin{enumerate}
\item Documentation of SE probe failure under a specific configuration (layer 25, TBG position, logistic regression, Llama-3-8B-Instruct, TruthfulQA), demonstrating that this configuration does not yield meaningful SE prediction.
\end{enumerate}

\begin{enumerate}
\item Identification of a replication gap between the experimental results and published benchmarks, suggesting that configuration details may significantly affect SE probe performance.
\end{enumerate}

\begin{enumerate}
\item A negative result that may inform future work on SE probing by highlighting the need for systematic configuration validation.
\end{enumerate}
